Cancerul de piele, în special melanomul, este foarte tratabil atunci când este
depistat precoce, însă șansele de supraviețuire scad în stadiile avansate; în
același timp, accesul la îngrijire dermatologică de specialitate este distribuit
inegal la nivel mondial. De aceea, scopul prezentei lucrări este de a prezenta un
model de inteligență artificială bazat pe învățare profundă, capabil să clasifice
leziunile cutanate pe baza imaginilor, precum și sistemul de ``screening al
cancerului de piele'' construit în jurul acestuia, accesibil atât într-un browser
web, cât și pe dispozitive mobile, care semnalează în ce măsură este justificată
examinarea de specialitate a unei leziuni cutanate.

Drept clasificator am folosit un model Vision Transformer (ViT) ajustat fin prin
învățare prin transfer, antrenat pe baza de date ISIC 2024, devenind astfel
potrivit pentru clasificarea imaginilor non-dermatoscopice de leziuni cutanate.
Am comparat modelul cu o rețea convoluțională ResNet-18 antrenată în condiții
identice, cu modelul EVA02 câștigător al competiției ISIC și cu modelul lingvistic
de mari dimensiuni Gemini 2.5 Flash. Pe setul de testare, modelul ViT ajustat fin
a atins un AUC de $0{,}939$, depășind substanțial modelul convoluțional de
referință și egalând nivelul EVA02.

Am extins modelul de clasificare antrenat într-o aplicație completă, cuprinzând
clienți mobili și web, un strat MCP care servește modelul, validarea imaginilor și
explicații în limbaj natural bazate pe Gemini, precum și stocare în cloud
autentificată și separată pe utilizatori.

\textbf{Cuvinte cheie}: cancer de piele, melanom, învățare profundă, Vision
Transformer, învățare prin transfer, ISIC 2024, model lingvistic de mari
dimensiuni, sănătate mobilă
